\section{Challenges and opportunities}
% \label{sec:challenges_and_opportunities}

% \textcolor{red}{
%   Government and Private Company Collaboration in the Governance of Shared Mobility Schemes: A Case Study of dockless Bike-Sharing Schemes in Sydney, Australia.
%   Nevertheless, DBSS are not perfect and their implementation presents significant challenges, particularly in terms of urban governance. Users sometimes dump or illegally park bikes in public spaces [13], disrupting traffic, creating hazards for pedestrians and affecting urban aesthetics [12]. In their initial response, some local governments in Australia impounded misplaced bikes, while the City of Melbourne went even further and announced a ban on DBSS in 2017 [14]. These fairly extreme early measures triggered fierce public criticism of government authorities for not supporting transport-sharing schemes in the public interest and for being averse to innovation. A more considered and sophisticated solution has emerged more recently in the form of collaborative governance (CG) between government actors and private DBSS companies [9,13,15]. CG is a term used to refer generally to cross-boundary collaboration between the public and private sectors in co-managing public programs or resources [16]. Although some scholars [10,17] have noted that the collaboration between multiple actors may make the decision-making process complicated, cross-sectoral collaboration can help establish reciprocal and trust-based relations and reduce conflict [18].
% }

% \textcolor{red}{
%   Despite the popularization of the bike-sharing concept during the last decade, the system design is not free of difficulties. The main operational problem is the system’s imbalance, caused by random and asymmetric demands (i.e. variable requests vs returns at different zones / stations). This gives rise to situations in which there are no available bicycles in a particular zone / station, while others are jam-packed. This last situation is more problematic in station-based systems because the user faces the impossibility of returning the bike in already full stations. System imbalance is addressed in the operating phase by artificially rebalancing bicycles among stations (or distribution zones, in case of free-floating systems). In addition to this, in the planning phase, the number, location and size of stations or distribution zones, and the size of the bicycle fleet, are critical decisions that affect rebalancing operations and steer the success of the implementation. Optimal system design, in terms of the selection of the previous decision variables, should respond to the optimization of the trade-off between users' and operating agencies' costs. On the one hand, from the users’ perspective, the service level is defined by the availability and accessibility of bicycles (i.e. in terms of the access distance), both at the origin and destination of the trip. This includes the ease of bicycle return. In general, the level of service increases with the number of bikes and with the number and size of stations. On the other hand, operating agencies aim to limit the investment and aintenance costs of bikes and stations, and the operational costs of bike repositioning within the system. The viability of a system depends on the ability to provide a reasonably good level of service while keeping costs affordable.
% }

%ChatGPT Thesis Review request (prinicipi conversa) (My Thesis project)
    % Data Gaps and Areas for Further Research

    % The intersection of public health and urban mobility that bike-sharing occupies is a relatively young field of research, and several important gaps remain. One such gap is understanding the long-term, population-level health impacts of bike-sharing. Many studies, like those cited above, have focused on immediate or annualized effects (e.g. deaths avoided per year, short-term injury rates). However, it is less clear how the proliferation of bike-sharing over a decade or more might influence epidemiological trends in a city. For instance, does sustained bike-share usage correlate with lower rates of obesity in the city’s population? Has it reduced healthcare costs related to sedentary lifestyle illnesses? These questions are not fully answered yet. The literature acknowledges that while individual-level benefits are well documented, few studies have assessed outcomes like community-wide changes in physical activity levels, long-term morbidity, or healthcare utilization due to bike-sharing
    % mdpi.com
    % . Mental health effects are another area that warrants deeper investigation. Anecdotal evidence and small-scale studies suggest positive mood and reduced stress from cycling, but no large-scale longitudinal study has yet quantified mental health improvements attributable to bike-share programs. Given rising interest in the mental well-being aspect of urban design, this could be a fruitful research avenue (for example, surveying bike-share users over time to measure changes in self-reported stress or happiness).

    % On the safety side, a critical data limitation is the under-reporting of cycling incidents, especially non-fatal and “near miss” events. Official crash statistics (police reports) often miss a significant portion of bicycle crashes – for example, if a cyclist falls or collides with a pedestrian and no motor vehicle is involved, it may never be recorded in police databases. Hospital records show many more cycling injuries than police records do, indicating that traditional data sources undercount the true number of incidents
    % itf-oecd.org
    % . This gap hinders our full understanding of bike-share safety. To address it, researchers are calling for more comprehensive data collection efforts: linking hospital and emergency room data with bike-share usage data, conducting regular user surveys about falls or close calls, and using new technologies (like mobile apps or crowdsourced reporting platforms) to capture minor crashes. Such data would help identify frequent causes of bike-share accidents (e.g. road surface issues, particular intersections of concern) that might not appear in coarse accident statistics.

    % Another emerging research direction is leveraging data-driven and AI methods to enhance our understanding of bike-share health and safety impacts. Since modern bike-sharing systems generate abundant data (GPS traces of trips, speed profiles, etc.), there is an opportunity to apply machine learning and predictive analytics to identify patterns associated with unsafe conditions. For instance, AI algorithms could analyze millions of trip trajectories to flag locations where users consistently slow down or divert – potentially revealing “near-miss hotspots” where cyclists feel unsafe. Likewise, integrating traffic data, pollution sensors, and health datasets could enable more sophisticated models that predict health outcomes under different scenarios (e.g. how much health benefit would doubling the bike-lane network yield, or how would a city-wide helmet campaign affect injury statistics). Such data-driven approaches are still in their infancy in this domain, but they align well with the thesis of applying AI in urban mobility analysis. Future research that combines transportation engineering, public health data, and AI could provide city planners with powerful tools to maximize the health benefits of bike-sharing while minimizing risks.

    % Finally, ongoing evaluation is needed as bike-sharing systems evolve. The rise of dockless bikes and e-bikes, the impacts of events like the COVID-19 pandemic (which in some cities led to surges in cycling), and changes in urban form all have potential health and safety ramifications. For example, during COVID-19 lockdowns, many cities observed a spike in cycling as people avoided public transit – understanding how bike-sharing contributed to safe mobility during that period is an active topic of study [81]. Additionally, researchers note that more work is required to quantify certain impacts precisely, such as the reduction in traffic congestion or air pollution attributable to bike-share (which indirectly affects safety and health). There is also room for natural experiments – for instance, comparing cities with similar profiles but different approaches to bike-share safety (one with mandatory helmets vs one without, or one with an extensive protected bike network vs one with minimal infrastructure) to rigorously evaluate which policies yield the best health/safety outcomes. In summary, the evidence to date overwhelmingly suggests that bike-sharing offers significant net health benefits and manageable safety risks. Increased physical activity has a positive impact that far exceeds the negative impacts of pollution exposure or crash injuries in most studied scenarios
    % pubmed.ncbi.nlm.nih.gov
    % . Nevertheless, to ensure these benefits are fully realized and equitably distributed, city authorities and researchers must continue to work on improving safety – through better infrastructure, education (e.g. promoting helmet use and safe riding practices), and inclusive planning – and to fill the knowledge gaps with robust data-driven research. This will enable bike-sharing systems to truly become a cornerstone of healthy, safe, and sustainable urban mobility.

% \textcolor{red}{
%   %Shaheen2010
% 3rd generation advances has helped detter thefts
% }

% Current limitations (e.g., vandalism, equity, fleet maintenance).
% "Convenience and safety" section -> Fishman2016

% Governance issues (private/public control, data ownership).

% Open challenges (infrastructure, regulation).

% Future opportunities (integration into MaaS, E-BSS expansion, digital twins).





% DeMaio2009 Bike-sharing: History, Impacts, Models of Provision, and Future
% Ricci2015 has a section called "Evidence on the impacts of BSS"


% Si2019 Mapping the bike sharing research published from 2010 to 2018: A scientometric review

% Zheng2020 The Development, Characteristics and Impact of Bike-Sharing Systems: A literature review
% Chen2020 dockless bike-sharing systems: what are the implications?
% Ma2020 Bike-sharing systems’ impact on modal shift: A case study in Delft, the Netherlands
